\section*{第 \ref{ch:b2:elements-in-the-earth} 章——地球中的元素}

\begin{solution}{exo:b2:elements-in-the-earth:1}
钾占 \ce{K2O} 质量的 $2\times39.10/(2\times39.10+16.00) = 78.20/94.20 = 0.8301$，
故 $\num{5130}\times0.8301 = \qty{4259}{ppm}$ 为钾。钛占
\ce{TiO2} 的 $47.87/79.87 = 0.5993$：$\num{27310}\times0.5993 = \qty{16370}{ppm}$ 为钛。
\end{solution}

\begin{solution}{exo:b2:elements-in-the-earth:2}
铝是\omterm{def:b2:elements-in-the-earth:goldschmidt}{亲石元素}且难熔：硅酸盐地球保留了全部铝，其比例接近
\omterm{def:b2:planet-formation-meteorites:chondrule}{球粒陨石}值（1.16）。钠是\omterm{def:b2:elements-in-the-earth:goldschmidt}{亲石元素}但易挥发（\qty{958}{K}），它在
吸积之前或吸积过程中已丢失，余下 0.22。镍与镁同时凝结
（\qty{1353}{K}），因此挥发性无法解释 0.077：镍是\omterm{def:b2:elements-in-the-earth:goldschmidt}{亲铁元素}，其中大部分
进入了地核。铱难熔（\qty{1603}{K}）且为强\omterm{def:b2:elements-in-the-earth:goldschmidt}{亲铁元素}：
0.0029，几乎全部在地核中。
\end{solution}

\begin{solution}{exo:b2:elements-in-the-earth:3}
$\mathrm{La}_N = 31/0.232 = 133.6$，$\mathrm{Yb}_N = 2.0/0.163 = 12.27$，于是
$(\mathrm{La}/\mathrm{Yb})_N = 10.9$。地壳富集所有\omterm{def:b2:elements-in-the-earth:ree}{稀土元素}，且
轻稀土的富集程度约为重稀土的十一倍：地壳由熔体构成，而熔体
最易容纳不相容的轻稀土。
\end{solution}

\begin{solution}{exo:b2:elements-in-the-earth:4}
$\mathrm{Sm}_N = 4.7/0.145 = 32.4$, $\mathrm{Eu}_N = 1.0/0.0546 = 18.3$,
$\mathrm{Gd}_N = 4.0/0.198 = 20.2$；$\mathrm{Eu}^* = \sqrt{32.4\times20.2} = 25.6$，并且
$\mathrm{Eu}/\mathrm{Eu}^* = 18.3/25.6 = 0.72$，为负异常。斜长石
能容纳 \ce{Eu^2+}，它从形成上地壳的熔体中结晶，并留存
于更深处。
\end{solution}

\begin{solution}{exo:b2:elements-in-the-earth:5}
$\bar D_{\mathrm{La}} = 0.65\times0.058+0.18\times0.001+0.17\times0.001 = 0.0377 +
0.0002 + 0.0002 = 0.038$；$\bar D_{\mathrm{Lu}} = 0.3445+1.044+0.0306 = 1.42$。镥是
\omterm{def:b2:elements-in-the-earth:partition}{相容元素}（$\bar D > 1$），镧则为强\omterm{def:b2:elements-in-the-earth:partition}{不相容元素}。石榴子石携带
$1.044/1.42 = 74$\% 的镥，尽管它仅占
岩石的 18\%。
\end{solution}

\begin{solution}{exo:b2:elements-in-the-earth:6}
对于钠，硅酸盐地球的 Na/Mg 为 $2670/\num{228000} = 0.01171$，而 CI
值为 $5010/\num{95870} = 0.05226$：比值为 0.224。对于镍，$1960/\num{228000}
= 0.008596$，对比 $\num{10640}/\num{95870} = 0.1110$：比值为 0.0775。钠
因易挥发而亏损（\qty{958}{K}）；镍在 \qty{1300}{K} 以上凝结，
其亏损程度是钠的三倍，因为它是\omterm{def:b2:elements-in-the-earth:goldschmidt}{亲铁元素}：地核
吸收了其中的大部分。
\end{solution}

\begin{solution}{exo:b2:elements-in-the-earth:7}
La：$5.21\times0.058 = \qty{0.302}{ppm}$；Lu：$0.53\times0.53 = \qty{0.281}{ppm}$。在
比值之比中，标准化值相互抵消：
$(\mathrm{La}/\mathrm{Lu})_{N,\mathrm{cpx}}/(\mathrm{La}/\mathrm{Lu})_{N,\mathrm{melt}} =
D_{\mathrm{La}}/D_{\mathrm{Lu}} = 0.058/0.53 = 0.11$：晶体的图案向重稀土方向
的斜率约为熔体的九倍。
\end{solution}

\begin{solution}{exo:b2:elements-in-the-earth:8}
地壳：$(12/0.658)/(31/0.648) = 18.24/47.84 = 0.38$。夏威夷玄武岩：
$(18.10/0.658)/(15.20/0.648) = 27.51/23.46 = 1.17$。MORB：$7.96/8.04 = 0.99$。只有
地壳带有明显的铌亏损，携带俯冲带的特征信号。
\end{solution}

\begin{solution}{exo:b2:elements-in-the-earth:9}
\omterm{def:b2:elements-in-the-earth:bse}{原始地幔}：$240/0.0203 = \num{11800}$。MORB：K $= 0.160\times0.830\times10^4 =
\qty{1328}{ppm}$，K/U $= \num{11200}$。地壳：K $= 2.80\times0.830\times10^4 =
\qty{23240}{ppm}$，K/U $= \num{8600}$。钾从地幔到地壳增加约一百倍，
而 K/U 的变化不到三分之一：两种元素的不相容程度大致相当，
熔融几乎不会将它们分开。铀难熔，因此其丰度由
\omterm{def:b2:planet-formation-meteorites:chondrule}{球粒陨石}模型给出；再由熔岩测得的 K/U 即可求出行星的钾含量，
而钾易挥发，不能取自\omterm{def:b2:planet-formation-meteorites:chondrule}{球粒陨石}。
\end{solution}

\begin{solution}{exo:b2:elements-in-the-earth:10}
在平衡状态下，对每种元素有 $C_{\mathrm{rock}} = \bar D\,C_{\mathrm{melt}}$（
熔体分数较小时岩石成分几乎不变），因此
$(C_1/C_2)_{\mathrm{melt}} = (C_{1}/\bar D_1)/(C_{2}/\bar D_2)$，此处 $C_1,
C_2$ 为岩石的浓度，该比值为 $\bar D_2/\bar D_1$。对于 La 与 Lu：$1.419/0.03805 = 37$。熔体
由含石榴子石的地幔生成，其中 La 相对于 Lu 的富集程度是源区的 37 倍：
石榴子石滞留了重稀土，熔体的图案陡峭。
\end{solution}

\begin{solution}{exo:b2:elements-in-the-earth:11}
对每种元素有 $C_{\mathrm{mix}} = fC_U + (1-f)C_M$，关于 $f$ 为线性。比值
$\bigl(f\,\mathrm{Th}_U + (1-f)\mathrm{Th}_M\bigr)/\bigl(f\,\mathrm{Nb}_U +
(1-f)\mathrm{Nb}_M\bigr)$ 为两个线性函数之商，即关于 $f$ 的双曲线。
当 Th/Nb $= 0.2$ 时，$f(10.5-2.4) + (1-f)(0.404-1.048) = 0$，即 $8.1f = 0.644(1-f)$，
$f = 0.644/8.744 = 0.074$：约 7\% 的地壳即可使 Th/Nb 由 0.077 升至 0.2。对于
$f = 0.1$: Th $= 1.05 + 0.364 = 1.414$, Nb $= 1.2 + 4.716 = 5.916$, Th/Nb $= 0.24$.
\end{solution}

\begin{solution}{exo:b2:elements-in-the-earth:12}
去除不含这两种元素的组分，两者的质量，
进而其比值保持不变。若铝全部被保留，$(\mathrm{Al}/\mathrm{Mg}) = 1.16$ 倍于 CI 值，
意味着硅酸盐地球保留了其\omterm{def:b2:planet-formation-meteorites:chondrule}{球粒陨石}镁的 $1/1.16 = 0.86$：
缺失 14\%。硅相对于铝为 CI 值的 $0.83/1.16 = 0.72$ 倍：
\omterm{def:b2:planet-formation-meteorites:chondrule}{球粒陨石}硅中将有 28\% 缺失。
\end{solution}

\begin{solution}{pb:b2:elements-in-the-earth:1}
\textbf{1.} Fe/Mg $= \num{182800}/\num{95870} = 1.907$; Ni/Mg $= 0.1110$; Si/Mg $=
1.111$; Al/Mg $= 0.0887$.

\textbf{2.} $(\num{182800}/55.85)/(\num{106500}/28.09) = 3273/3791 = 0.863$，
即每个硅原子对应 0.863 个铁原子，与光球的 0.81 相差在 7\% 以内：这是合理的，因为 \omterm{def:b2:elements-in-the-earth:ci}{CI 球粒陨石}
在造岩元素上与太阳相符。

\textbf{3.} 去除不含比值中任一元素的组分，
该比值保持不变（见\cref{exo:b2:elements-in-the-earth:12}）：只有绝对浓度发生变化。

\textbf{4.} 镁为\omterm{def:b2:elements-in-the-earth:goldschmidt}{亲石元素}，不进入地核；它是主量元素，
在\omterm{def:b2:planet-formation-meteorites:chondrule}{球粒陨石}与地幔中均测定得很准确。铁是在地核与地幔之间分配的元素：
它不能充当自身的参照。

\textbf{5.} 一起凝结的元素在星云中未因挥发性而分离，
因此整个地球的比值应保持不变。Al/Mg 出现偏离：铝为
\omterm{def:b2:elements-in-the-earth:volatility}{难熔元素}（\qty{1653}{K}），硅酸盐地球的 Al/Mg 是 CI 值的 1.16 倍。

\textbf{6.} Fe $= 8.05\times55.85/71.85 = 6.26$\%；Ni $= 0.25\times58.69/74.69 =
0.196$\%。

\textbf{7.} Fe/Mg $= 6.26/22.8 = 0.2746$；Ni/Mg $= 0.196/22.8 = 0.00860$。

\textbf{8.} 铁：$1 - 0.2746/1.907 = 0.856$，缺失 85.6\%；镍：
$1 - 0.00860/0.1110 = 0.923$，缺失 92.3\%。

\textbf{9.} 铱的亲铁性强于镍，镍又强于铁：
元素越亲铁，金属相带入地核的该元素就越多。

\textbf{10.} 在硅酸盐中，铁以 \ce{Fe^2+} 形式存在于橄榄石和辉石中（FeO 为 8\%）：
被氧化为 FeO 的铁与氧结合，行为如\omterm{def:b2:elements-in-the-earth:goldschmidt}{亲石元素}，留在
地幔中；只有金属铁下沉。

\textbf{11.} $(1-x)C_{\mathrm{Fe}} + xc_{\mathrm{Fe}} = r_{\mathrm{Fe}}(1-x)C_{\mathrm{Mg}}$;
除以 $(1-x)c_{\mathrm{Fe}}$ 即得命题。

\textbf{12.} $r_{\mathrm{Fe}}C_{\mathrm{Mg}} - C_{\mathrm{Fe}} = 1.907\times22.8 - 6.26 =
37.22$; $y = 37.22/85.5 = 0.4353$, $x = y/(1+y) = 0.303$.

\textbf{13.} $0.1110\times22.8 - 0.196 = 2.335$; $y = 2.335/5.2 = 0.4490$, $x = 0.310$.

\textbf{14.} $y = (85.5\times37.22 + 5.2\times2.335)/(85.5^2 + 5.2^2) =
3194.4/7337.3 = 0.4354$, $x = 0.303$.

\textbf{15.} 每个方程的权重取其系数的平方：$85.5^2/5.2^2 = 270$
倍，对丰度高得多的铁而言。

\textbf{16.} $37.22/2.335 = 15.9$，与模型中的 $85.5/5.2 = 16.4$ 相比，
相差 3\%。

\textbf{17.} $0.303\times\num{5.972e24} = \qty{1.81e24}{kg}$.

\textbf{18.} 比 32.5 低 $(32.5-30.3)/32.5 = 7$\%。

\textbf{19.} $y = 0.325/0.675 = 0.4815$，地核中铁 $= 37.22/0.4815 = 77.3$\%，
镍 $= 2.335/0.4815 = 4.85$\%：轻元素占 18\%，约为密度所允许的 5\% 至
10\% 的两倍。地球的 Fe/Mg 应略高于 CI 值，
或其硅酸盐中的镁高于模型值。

\textbf{20.} 硅、硫、碳、氧和氢。硅在
硅酸盐地球中含量偏低：其以镁归一化的比值 0.83，低于\omterm{def:b2:elements-in-the-earth:volatility}{难熔元素}的水平。

\textbf{21.} $1.111\times22.8 - 21.0 = 4.331$；$4.331/0.4354 = 9.9$\% 的硅位于
地核中。

\textbf{22.} 9.9\% 对比 6\%：地核中的硅可能多于模型所给，或
部分硅从未被吸积。硅在 \qty{1310}{K} 时凝结，略
低于镁和铁（1336 与 \qty{1334}{K}）：部分硅可能留在了气体中。

\textbf{23.} 地核：$0.325\times0.5 = \qty{0.1625}{ppm}$（占地球质量）；硅酸盐
地球：$0.675\times0.001 = \qty{0.000675}{ppm}$。地核所含
金的 $0.1625/0.1632 = 99.6$\%。

\textbf{24.} 仅凭化学即可得出
地核占地球质量的 \textbf{\qty[text-series-to-math=true]{30.3}{\percent}}，
与地震学测得的 32.5\% 相差在 7\% 以内。
\end{solution}
